\section{Discussion}\label{sec:discussion}

\subsection*{Main findings and interpretation}

When the EEG carried no label information, covariate adjustment created discrimination of up to 11.9 pp in nonlinear classifiers on real EEG features and at most 1.0 pp in logistic regression. The same pattern appeared in an independent cohort recorded on a consumer device with different covariates (12.7 and 11.6 pp in the tree classifiers; at most 0.3 pp in logistic regression), so the effect is not specific to one archive. A reference fitted on training controls leaked even with a linear classifier (null AUC 56.4).
The decision rules that turned these null experiments into verdicts rarely raised false alarms (R1: 0.1\% of draws in simulation cells without leakage) and detected leaks of known size (88\% for V0 with 20 replicates, 96\% for simulation cells with a long-run AUC above 55); between 52 and 55 they were mostly ``uncertain''. Their weak point is the paired rule for high-variance pipelines, where 20--30 replicates often did not reach a decision and declared leakage where none was possible; with the pools extended to 120 replicates, the mean differences shrank from +2.2--3.0 to +0.9--1.1 pp and false alarms fell to 2.7\%--4.2\% with 50 replicates, which points to too few replicates rather than a defect of the rule.
Rank-normal before adjustment reduced the leak in tree classifiers but not enough in the Bi-LSTM, and its ``no evidence'' decisions changed with thresholds 1 pp stricter or looser, so the mitigation is reported as conditional on the representation, the classifier, and the rule. In the simulation with a fixed marginal distribution, the leak of window-level adjustment was not shown to grow with the correlation of the windows within patients (change $-$1.8 pp from ICC 0 to 0.95; Holm 0.79), and the advantage of aggregating to patient means before adjusting (+8.9 pp at ICC 0 in one simulated cell) vanished at ICC 0.95 (conjecture: patient means keep the tails when the windows are strongly correlated).
For the TUEP label, F2 (69.7; 95\% CI 61.0--78.5), I2 (68.0), and C1 without adjustment (67.5) could not be distinguished from one another or from the metadata baseline (65.5; $p = 0.45$--$0.84$), and, conditional on the stored candidates, there is no evidence that F2 exceeds the metadata baseline or C1 (corrected AUC 67.4 and 65.6; max-$T$ $p \geq 0.36$). F2 is therefore an audited test case, not a validated final pipeline; it contains NM, which did not fully pass the audit, and the linear variant without NM (68.2) leads to the same conclusions. The similarity of F2, I2, and C1 may reflect the resolution of this design (200 patients; 0.86 events per predictor in I2; smallest detectable differences of 7 to 9 pp) as much as similar information content.
Adding the EEG features to all non-EEG predictors raised the per-fold AUC by 7.6 pp, but the corrected test did not establish this gain ($p = 0.082$; Holm 0.14), and refitting the label-free transforms inside the inner folds changed the AUCs by at most 0.3 pp for every model, including NM and F2.

Discrimination of the TUEP label persisted among case--control pairs matched on year, device, and age within 5--15 years (F2 72.0--73.0; I2 and C1 similar, so the signal was not shown to be specific to F2), with equal weight per stratum (60.4--73.6) and without the largest stratum (73.4--75.5). Matching restricts the comparison to 141--173 patients who have a partner and does not remove confounding by medication, report-derived labels, or unmeasured acquisition factors \cite{maxion2024}; the persisting score is therefore not a disease marker. This is evidence that the labels carry signal beyond age, year, and device, not evidence of robustness across populations or of clinical applicability.

\subsection*{Comparison with the literature}

The leak of per-feature adjustment into nonlinear classifiers is the mechanism described by Hamdan et al. \cite{hamdan2023}; this study does not claim it as new. What is added is its size on real EEG features from two archives with labels that carry no EEG information, the error frequencies of explicit rules for the feature-shuffle test on clustered data, and the observation that the mitigation that reduced the leak in trees did not do so enough in a sequence model. The control-referenced adjustment probably leaked through a fitting asymmetry: only the residuals of the training controls are orthogonal to the design ($3.9\times10^{-15}$; residual of the training patients with epilepsy 1.39), so the coefficient error carries age, sex, and era to all patients (conjecture). No tail property that explains the leak was identified: skewness, residual quantiles, and the symlog kink did not explain it better than kurtosis, so the map is empirical. In neuroimaging, covariate regression outside the folds lowered performance \cite{rosenblatt2024}, so the direction of the effect can depend on the design.
Against CBraMod-FT, F2 was higher in the patient-level bootstrap but not in the corrected test (limited and exploratory evidence). In the same pipeline, frozen CBraMod (66.2) and LaBraM (63.6) did not differ significantly from I2, BENDR (53.2) was lower (Holm 0.03), and the 50-feature anatomical compression (58.8) did not raise the AUC. Pretraining on TUH, which contains TUEP, may favor foundation models by an unmeasured amount \cite{jiang2024,wang2025cbramod}.
Reproduction of published methods under the patient-level protocol gave a pooled AUC of 54.9 (Tang-F) and 55.0 (Rohira-F), far below the accuracy above 95\% under random splits \cite{rohira2023,tang2024}; accuracy and AUC are different metrics, and the gap cannot be attributed to patient leakage alone because implementations also differ.

For managers of clinical EEG archives, this study suggests five checks before performance figures from stored recordings are interpreted as disease information. (i) Split the data by patient before using any data statistics \cite{delpup2025,brookshire2024,rosenblatt2024}. (ii) Report an acquisition-metadata baseline and a model with all non-EEG predictors. (iii) Evaluate within year $\times$ device strata with an age caliper and report who contributes to the pairs. (iv) Run the feature-shuffle test with and without adjustment, with enough replicates for the variance of the pipeline (at least 50 when the SD of the paired difference is near 8 pp; Table~\ref{tab:oc}), especially when adjustment is combined with nonlinear classifiers. (v) Report selection optimism over the stored candidates when a method is chosen after many rounds on the same patients. These checks are evaluation recommendations, not a clinical standard, and the results do not show clinical readiness.

\subsection*{Strengths and limitations}

The strengths of this study are: patient-level splitting with 50 canonical folds uniform for all models; comparators under the same protocol (published methods reproduced per patient, foundation models under equal treatment); non-EEG baselines, EEG beyond all non-EEG predictors, and analysis within strata and age calipers; null and positive controls for the audit rules; and version-controlled analysis plans with independent check groups for rounds 39--50.

Data source: the case study comes from one retrospective archive ($n = 200$; first sessions only) without temporal validation; the plasmode was repeated in one independent public cohort with a different device, population, and covariates (Emotiv), but the case-study results (F2, baselines, strata) were not replicated in an independent clinical archive, and no result supports clinical applicability. Transfer of predictions to two other cohorts that differ in task, population, and device was unsuccessful (Kural 57.1; Emotiv 56.5 with frozen statistics); these transfer analyses do not test the audit. Labels derived from reports and clinical history (medication and EEG) may be circular with respect to the signal, and the direction of the bias cannot be determined \cite{veloso2017}. The non-epilepsy spectrum cannot be characterized (TUEP v2.0.1 has no clinical reports); the results do not apply to first seizures or later sessions.
Analysis history: 50 rounds (44 design rounds and six follow-ups, 45--50) and one post hoc analysis were run on the same patients and folds without a hidden test set. Every plan was written after earlier results were known, so all inferences are exploratory with respect to this history, and Holm within a round does not control error across rounds. The selection optimism of 2.0 to 3.8 pp is conditional on the stored candidates and does not cover the other design degrees of freedom. The wide CIs and the post hoc equivalence test (no pair shown to be equivalent) mean that the absence of a significant difference does not equal equivalence. Differences below about 7 to 9 pp between F2, I2, and C1 (about 16 pp against metadata) cannot be ruled out.
Audit: the rules were calibrated on stored replicates of this archive and on simulations, and for high-variance pipelines the paired rule needs 50 or more replicates and still returns ``no evidence'' in only about half of the draws when the true difference lies near 1 pp; the $\gamma=0$ control failed for XGBoost (7 of 10 cells), NM-V4 did not fully pass (+1.9 pp), and decisions about the mitigations depended on the thresholds. The thresholds (52.0, 54.0, and 55 for R1; 1.0 and 3.0 pp for R2) are operational conventions fixed before the analysis of their properties (round 49), not derived from theory. The SE in the rules reflects Monte Carlo error on the same patients, not how the size of a leak would change in a cohort with a different label--covariate association.
Comparators and scope: published comparators and foundation models were evaluated with official recipes without tuning; the simulation and plasmode map leakage, not diagnostic performance. Only rank-normal, log, and symlog transforms, aggregation to patient means, inverse-propensity weighting, and linear classifiers were evaluated as mitigations (cross-fitting only on raw features); other harmonization strategies were not, so the conclusions on mitigation apply to these. Calibration and decision curves at a sample prevalence of 50\% without a defined clinical decision are descriptive only. Future research needs independent archives with acquisition metadata and patient blocks in which the frozen audit can be replicated, medication data, and prospective evaluation with neurologists.
